\section{What The Chapter Carries Forward}

The chapter has kept the single invariant and made it charged. The electron
was read as the \(-1\) value of the discriminating action over the flat
baseline. The positron was read as the \(+1\) value of the same pair grammar
only after the baseline tilted. The two signs did not arrive as free
opposites. They were forced by the relation between second variation,
baseline, and discriminated residue.

The chapter also kept identity narrow. It used the local quotient that lets
a selected baseline stand as the same baseline for a specified reading, and
it used the motion grant that lets stationarity answer to the world. It did
not add a third source of identity. Under those two grants, boundary geometry
generated the gauge by producing the action that reads the charged pair.
The gauge did not replace geometry, and geometry did not swallow the gauge.
Their finite carrier split.

That split is what the next chapter receives. Since the geometric and gauge
faces do not form an interior direct sum, their reconciliation cannot be an
interior scalar mixed term. The grammar must look to the boundary. The
zero-mesh Cauchy reading and the finite Cohen forcing ledger have supplied
the continuum discipline needed for that next move. What remains is to read
how the boundary carries the residue that the interior is forbidden to mix.

Chapter 05 begins there. The present chapter has forced the charged pair,
the baseline dependence, the geometry-to-gauge conversion, the split, and
the finite forcing of the continuum reading. Boundary radiation is not an
extra topic after those results. It is the next grammatical obligation.

